\subsection{The Joint Episode as Kairos}

In the metaphysics of becoming, the present is not an instantaneous point but a structured temporal window whose internal organization conditions what can emerge next. The Joint Episode realizes one such coherent, empirically detectable \textit{kairos}: a bounded interval whose joint satisfaction of local persistence intensification (Layer 1) and relational scale alignment (Layer 2) materially raises the conditions of possibility for qualitative systemic change.

As illustrated in Figure~\ref{fig:three_layers_je}, the kairos is not merely a period of low antisynchrony, nor simply a burst of local trapping. It is the specific conjunction in which a maximal interval of low $A(k)$ contains at least one significant excursion of the critical-mass metric $M$. This conjunction defines a window during which the system is simultaneously \emph{intensified in its local becoming} and \emph{permeable across scales}. Within such a window, local structure is no longer isolated; it participates in a temporarily unified field of persistence whose relational coherence may enable or amplify transitions that would otherwise remain improbable.

The notion of kairos therefore carries both a descriptive and an ontological weight. Descriptively, it names the observable temporal unit recovered by the detection algorithm. Ontologically, it designates the form of the present in which the potential for reorganization is not abstract but locally and relationally enacted.

\subsection{Act of Being and Futurization}

Drawing on the three-layer framework, the Joint Episode is the locus at which the act of being---understood as the local intensification and maintenance of persistence---acquires the additional relational conditions required to participate in global becoming. The act is not exhausted by Layer 1 excursions of $\tau_s$ and $M$; it is completed, in the relevant sense, only when those excursions are embedded within a sustained low-$A(k)$ interval.

This embedding makes possible what may be termed the \textit{futurization of the present}. In ordinary succession, the present passes into a future that remains largely external to it. During a Joint Episode, however, the present is internally reconfigured: its local persistence patterns are brought into sufficient scale coherence that subsequent structural reorganization can be understood, at least in part, as an outgrowth of the episode's own organization rather than as an entirely exogenous event. The system does not merely ``have'' a future; through the kairos it begins to enact configurations that prefigure or constrain that future.

Importantly, this futurization remains probabilistic and non-deterministic. The synthetic results (Section 6) show that even under idealized labeling the discrimination between episodes that precede detectable Layer-3 events and those that do not is modest. The Joint Episode supplies necessary relational conditions; it does not furnish a sufficient causal bridge. This limitation is itself philosophically instructive: it indicates that the ``act of being'' realized in the episode is enabling rather than determining.

\subsection{Non-Reductive Hierarchy}

The three-layer ontology is deliberately non-reductive. Layer 1 (local intensification) is not ontologically prior to Layer 2 (relational coherence) in a way that would allow the latter to be derived from the former by aggregation; nor does Layer 3 (structural reorganization) stand as a higher-level effect that could be explained by summing independent Layer-1 events. Each layer possesses its own primitive descriptors and its own form of irreducibility.

The Joint Episode respects this irreducibility while identifying their joint realization as a new primitive unit. The unit is not reducible to Layer 1 alone (a mere local burst of persistence would lack the permeability required for propagation), nor to Layer 2 alone (sustained low antisynchrony without intensification lacks the local ``mass'' that could seed reorganization). Its identity is irreducibly relational-temporal: it exists only at the intersection realized over a maximal interval.

This non-reductive stance has methodological consequences. Detection algorithms must track both layers simultaneously; explanatory claims must be careful not to collapse one layer into another; and theoretical development must treat the Joint Episode as an ontologically basic category rather than an epiphenomenon of either local dynamics or global statistics.

\subsection{Implications from Synthetic Validation}

The controlled experiments reported in Section 6 provide an important reality check on these ontological claims. Even when episodes were labeled using an oracle that possessed perfect knowledge of the injected structural perturbations, the best scalar score $J_{0.7}$ achieved only a modest AUC of approximately 0.476, with separation between pre-transition and non-transition episodes that was statistically indistinguishable from chance under the Mann--Whitney test. This result does not invalidate the framework; it clarifies its current scope.

In low-dimensional, few-module systems such as the four-oscillator Lorenz generator, the dominant statistical signature recovered by the Joint Score appears to be the L1--L2 coherence itself rather than a strong anticipatory marker of Layer 3 reorganization. The honest diagnostic reading is therefore that the Joint Episode, as currently operationalized, functions reliably as a detector of ontologically meaningful temporal units but does not yet provide a powerful predictor of structural transitions in minimal synthetic settings.

This finding constrains the ontological interpretation without negating it. It suggests that the ``potential for reorganization'' embodied in the kairos may remain latent until the system possesses either higher internal complexity, stronger coupling between scales, or additional mechanisms (for example, slow variables or memory across episodes) that amplify the relational conditions into actual reorganization. Future work on richer generators and empirical multivariate series will be required to determine whether the modest oracle ceiling observed here is an artifact of the testbed or a deeper feature of the present framework.

The synthetic results thus serve a dual purpose: they validate the operational detectability of the proposed unit while simultaneously exposing the limits of strong anticipation claims in the systems examined to date. Such honesty is essential if the ontological reading is to remain grounded rather than aspirational.